\section{在缝合曲面上移动插入}\label{subsec:motion}

布设构造将原来各表表示在每个单元中的一对球面上。现在我们把这些表合并到少数标记点处，只需利用洛朗展开的两类结果。在插入标记点处，取一个系数，可将移动中的表与原来就在该处的表融合。在接缝处，取一个阶数足够小的负阶系数，可将移动中的表移至相邻单元，并降低接缝能量。我们利用这一下降量，既能限制穿越次数，也能限制任一单元所能接收的能量。

在每个单元中，将原不可移动站点的标记点以及另一个真空标记点指定为\emph{汇集点}。它们留在各自的单元内。依次处理原来的可移动站点，跳过已经被吸收的站点。移动中的表保留其普通比特端口。若该表现有的两条腿在缝合球面 $A$ 和私有球面 $B$ 上的能量分别为 $m_A,m_B$，令 \[
       m=m_A+m_B/P,\qquad \mu=m+H,\qquad L=\lceil 2P\mu\rceil .
\] 其中 $P$ 是式~\eqref{eq:plant-scales} 的换算率。从一次移动开始直至结束，储备能量 $H$ 始终随这次移动同行，而 $m$ 则在每次融合或穿越后更新。

\subsection{由碰撞重建取值}

固定外层求和的一个分支，为推导下述恒等式也暂时固定私有选择器。分支本身就决定了能量，因此同一恒等式适用于所有容许的选择器取值。固定所有其他位置和插入表。收缩所得的每一表项都是移动者的 $B$ 坐标 $z$ 的有理函数，其 $A$ 坐标为 $z^P$。球面核只可能在当前球面上的其他标记点处有极点，且在无穷远处有界。在 $B$ 坐标为 $z_0$ 的标记点附近，令 $s=z-z_0$。由于 $P$ 是特征的一个幂，$A$ 上相应的坐标差为 $s^P$。因此，阶数为 $i$ 的 $A$ 融合与阶数为 $v$ 的 $B$ 融合对收缩的 $j=Pi+v$ 阶项作出贡献。在加权能量为 $m_0$ 的插入处，合并所得的能量为
\begin{equation}
            m_{\rm new}=m_0+m+j/P.                     \label{eq:motion-fusion}
\end{equation} 事实上，由命题~\ref{prop:fusion}，每次融合都将其阶数加到参与融合的两个能量之和上。两个输出能量非负，故 $i$ 与 $v$ 均有下界。因此极点阶数有界，对每个 $j$，只有有限多个数对 $(i,v)$ 作出贡献。

为重建当前位置的取值，使用单元中其他标记点处的下列系数。\begin{itemize}
\item 对每个插入，取其整个主部。\item 在指定的真空汇集点，还取 $0,\ldots,gL$ 阶各项，即使该汇集点已经吸收了一张表，仍然这样做。\item 在每条接缝处，限于所给的极点阶数界，仅取 $j\le -L$ 的各项。\end{itemize}
若函数的所有这些系数都为零，则它的极点只可能出现在至多 $g$ 条接缝处，且每处的极点阶数至多为 $L-1$。它在指定汇集点的零点阶数至少为 $gL+1$，在无穷远处没有极点。射影直线上非零有理函数的除子次数为零，故此函数恒为零。因此，在具有规定极点界的函数空间上，于移动者当前位置取值的泛函是上述系数泛函的线性组合。此组合逐表项作用，并且与各张表的具体内容无关。

插入处的一个系数是对所有相容数对 $(i,v)$ 的求和。对每一数对，施行两个融合映射，并将它们与移动表及静止表收缩，计入所选次序的符号。与汇集点碰撞便结束本次移动。在任何其他插入处只使用 $j<0$，因此合并后的表继续移动时所带的加权能量至多为原先两者之和。这一步可以吸收一个辅助插入、一个尚未轮到的原站点，或一个由先前穿越所留下的私有插入。

\subsection{穿越接缝}

设在能量为 $e$ 的一条接缝处选取了 $j\le-L$ 阶的系数。只有球面 $A$ 在此处缝合。接缝处 $B$ 上的插入是真空，因此 $B$ 展开的泰勒阶数为 $v\ge0$。对这样的每个 $v$，若 $i=(j-v)/P$ 为整数且落在融合范围内，则接缝处新的 $A$ 能量满足
\begin{equation}
       e'=e+m_A+i\le e+m_A-2\mu\le e-\mu .             \label{eq:motion-drop}
\end{equation} 满足 $e'<0$ 的项均为零。固定 $i$，用此前未使用的正则位置上的取值对 $B$ 的泰勒系数进行插值。在所得每一项中，近侧留下一个能量为 $m_B$ 的私有插入，其在 $A$ 上为真空。一条普通的相等连线将它的指标传至原表。我们称这一私有插入为本次穿越的\emph{残留插入}。

为了确定新的移动表，由线性性可以固定其原来的 $A$ 态 $x$。令 $a,b$ 分别遍历能量 $e,e'$ 的楔积基矢，并以 $C_{\rm near},C_{\rm far}$ 表示余下的两个球面形式，其中略去未参与本步变换的变元。引理~\ref{lem:seam} 给出恒等式 \[
\begin{aligned}
 &\sum_b C_{\rm near}(b)\sum_a [I_{e'}(a,x)]_b\,C_{\rm far}(a)\\
 &\hspace{1em}=\alpha\sum_{t\le m_A}\sum_b C_{\rm near}(b)
                      \sum_a [I_e(b,S_t x)]_a\,C_{\rm far}(a).
\end{aligned}
\] 在右端，内层求和是在远侧 $A$ 球面上插入 $S_t x$ 趋近 $b$ 时的 $e-e'-t$ 阶洛朗系数。该阶数为整数，因为非零项中的 $e'-e-m_A=i$ 与 $m_A-t$ 都是整数。用远侧正则位置上的取值对此系数进行插值，在相应的 $B$ 位置取真空。对 $b$ 求和就在能量 $e'$ 处将两球面缝合。新的移动表现位于远侧，能量为 $(t,0)$；它保留普通端口以及通向残留插入的连线（图~\ref{fig:motion-crossing}）。所有重排符号都只依赖指定的能量，即使 $x$ 与私有指标或某个未参与变换的变元相耦合，恒等式依然成立。

\begin{figure}[!htbp] \centering
\begin{tikzpicture}[x=1cm,y=1cm, font=\small,
 cell/.style={draw=ink!50,fill=ink!3,rounded corners=2pt},
 tab/.style={draw=ink,fill=white,rounded corners=1pt,inner sep=3.2pt,
             font=\small,align=center}]
 \node[font=\small\bfseries] at (-3.62,1.45) {近侧单元};
 \node[font=\small\bfseries] at (3.62,1.45) {远侧单元};
 \draw[cell] (-5.36,-1.11) rectangle (-1.88,1.13);
 \draw[cell] ( 1.88,-1.11) rectangle ( 5.36,1.13);
 \node[align=center] at (-3.62,.55) {缝合球面 $A$};
 \node[align=center] at ( 3.62,.55) {缝合球面 $A$};
 \fill[sea] (-1.88,.55) circle (2pt);
 \fill[sea] ( 1.88,.55) circle (2pt);
 \draw[sea,line width=.9pt] (-1.88,.55) -- (1.88,.55);
 \node[fill=white,inner sep=2pt,text=sea] at (0,.55) {接缝能量 $e'$};
 \node[tab] (res) at (-3.62,-.40) {私有残留插入\\ $(0,m_B)$};
 \node[tab] (mov) at ( 3.62,-.40) {移动表\\ $(t,0)$};
 \draw[dashed,rust,line width=.75pt] (res.east)--(mov.west);
 \node[fill=white,inner sep=2pt,text=rust] at (0,-.40) {$B$ 态指标};
\end{tikzpicture} \caption{能量为 $(m_A,m_B)$ 的移动插入从近侧单元穿越到远侧单元。实线表示将能量为 $e'\le e-\mu$ 的接缝状态彼此等同。虚线把原来的私有态传给新的移动表，该表在缝合球面上的能量为 $t\le m_A$。} \label{fig:motion-crossing}
\end{figure}

两次插值均可在原收缩所在的同一个有限域 $K$ 上完成。乘以一个公分母，就能将极点受指定限制且在无穷远处有界的有理函数空间嵌入次数不超过总极点阶数上界的多项式空间。在多于此数的正则域点上取值给出单射，所以这些取值泛函张成对偶空间，并能恢复各个所需的洛朗系数。排除所有已被占用的位置。映射 $z\mapsto z^P$ 是 $K$ 上的一个置换，因此可以在两个球面上使用对应位置作正则取值。下面的估计将次数控制在 $O_D(T)$，将每个单元内占用的位置数控制在 $D$ 的某个多项式以内。我们可以先在代数闭包上取得这些界，再据以在进行任何插值前选定 $K$。固定域元素的一种排列顺序。对依该顺序排列的每组容许正则点，在表字典中存入某个重建矩阵的各表项。局部控制只需指定所选点并检查它是否可用。

每次移动至多允许 $U=D^5$ 次穿越。达到此限时，移动者当前的标记点成为汇集点。在相邻两次穿越之间，每次非终止的融合都会消去一个静止插入。一次移动结束后不会再放出新的移动者，因此对每个求和项，该过程都将终止。

\subsection{移动的能量与连线估计}

先估计最终各插入存储的能量。元数为 $t_i$ 的原站点的初始加权能量为 $O((1+t_i/T)^2H)$。辅助插入在每个被占用的单元至多贡献 $DH$。另外，给每个原可移动站点分配一份储备 $H$。在每个单元中，初始插入能量和储备能量之和为 $O(D^2H)$。计入该批次的全部网络，对其所有单元求和所得的界为 \[
 O\!\left(H\left(D\,\#\{\text{被占用的单元}\}
                  +\#\{i\}+T^{-2}\sum_i t_i^2\right)\right)
       =O(WH/T).
\]

由式~\eqref{eq:motion-fusion}，可继续移动的融合不会增加参与融合的能量之和。只保留当前移动的储备。在一次穿越中，能量 $m$ 足以支付远侧的 $t\le m_A$ 和近侧残留插入的 $m_B/P$。在终止移动的融合中，汇集点可能获得的能量至多为 $m+gL/P\le(1+3g)\mu$，因为 $H\ge1$ 且 $P\ge1$。对汇集点的能量赋予权重 $1/(1+3g)$，对其他能量（包括储备）赋予权重一。在上述三种操作中以及移动达到穿越次数上限时，加权总和均不会增加。因此在所得每个求和项中，插入的总能量为 $O_g(WH/T)$。

每个单元内也有相应的界。一次穿越带入的能量计费至多为 $t+H\le\mu$，而由式~\eqref{eq:motion-drop}，接缝损失的能量至少有这么多。一个单元至多有 $g$ 条接缝，每条接缝的初始能量为 $O_g(D^2H)$。再加上该单元的初始计费量，就得到每一步中该单元当前插入能量的上界 $O_g(D^2H)$。此外，每次穿越都使接缝能量至少减少 $H$，因此在整个构造中，与一个单元关联的穿越共发生 $O_g(D^2)$ 次。只有这些穿越可以引入新的静止残留插入。由于 $H=O_D(T)$ 且 $P$ 固定，这些能量界与标记点个数界就给出了所述的总极点阶数上界 $O_D(T)$。在所有单元中，达到次数上限而结束的移动次数至多为穿越总次数除以 $U$。计入最初的不可移动站点与指定的汇集点，汇集点总数为 $O(W/(DT))$。

还须布设展开产生的普通连线。由式~\eqref{eq:plant-widths}，每次穿越留下的私有态占 $O_g(DT/\sqrt F+1)$ 比特。一端跟随移动者，另一端停留在残留插入处，直至被另一插入吸收，再随其移动。仍可能被吸收的残留插入只携带其自身那次穿越所产生的端点。移动端点一旦完成一次移动，其汇集点便固定不动。重复访问也计数时，一次移动至多访问 $U+1$ 个单元，每次访问至多收集 $O_g(D^2)$ 个残留端点，且整次移动至多产生 $U$ 个自身的端点。因此在任意一步跳移中，它所携带的端点数至多为 $O_g((U+1)D^2)$。利用与单元相关联的穿越次数界，并将终止情况计入，可得稍宽的界，每个单元有 $O_g((U+1)D^6)$ 段局部连线。单一单元内连续的移动只需一条局部连接。对第~\ref{subsec:planting} 节所取的 $F$，这些线段在该单元的总宽度为 $O(D^{-2\ell-1}T)+O_D(1)$。

私有选择器的连线从读取选择器的标记点，以及校验选择器和计算球面形式的站点出发。其端点沿相同的路径移动。每个单元初始拥有的此类端口数至多为 $D$ 和组数的某个多项式，因此由路径长度与穿越次数界可知，与一个单元关联的选择器线段及其副本的数量也受同一组参数的多项式控制。该计数与 $b$ 和选择器的编码率无关。固定组数后，选择足够低的编码率，使每个单元中的选择器总宽度（包括执行相等复制所需的宽度）至多为 $D^{-2\ell}T$。下一节将把剩余能级换算为各表的元数。
